% !TEX encoding = UTF-8
\documentclass[12pt,a4paper]{article}

% ===== Encoding & Language =====
\usepackage[utf8]{inputenc}
\usepackage[T1]{fontenc}
\usepackage{lmodern}

% ===== Math =====
\usepackage{amsmath, amssymb, amsfonts}

% ===== Graphics =====
\usepackage{graphicx}
\usepackage{caption}
\usepackage{subcaption}

% ===== Tables =====
\usepackage{booktabs}

% ===== Layout =====
\usepackage{geometry}
\geometry{margin=1in}

% ===== Links =====
\usepackage{hyperref}

% ===== Bibliography =====
\usepackage[numbers,sort&compress]{natbib}

% ===== Optional (recommended for better formatting) =====
\usepackage{setspace}
\usepackage{titlesec}


\title{\textbf{Evaluating Retrieval-Augmented Generation for Medical Question Answering and Literature Summarization}}

\author{
\\
\textit{Estin Bejaia} 
}

\date{\today}

\begin{document}

\maketitle

\begin{abstract}
Large Language Models (LLMs) have demonstrated strong capabilities in natural language understanding and generation, leading to their increasing use in medical question answering and literature summarization. However, a critical limitation of these models is their tendency to generate hallucinated or unverified information, which can be particularly harmful in healthcare contexts. This issue is amplified by the growing reliance of individuals on LLMs for medical advice, sometimes as a substitute for professional consultation.

In this work, we investigate the effectiveness of Retrieval-Augmented Generation (RAG) techniques in reducing hallucinations and improving factual correctness in medical applications. We propose a Graph-based RAG approach that leverages structured relationships and optimized querying to provide reliable contextual information to the LLM. This enables the generation of more accurate and evidence-supported responses.

While this work is still exploratory and does not yet define a novel contribution, it aims to evaluate the potential of graph-enhanced retrieval mechanisms in improving the safety and trustworthiness of LLM-based medical systems.
\end{abstract}

\section{Introduction}

The rapid advancement of Large Language Models (LLMs) has significantly impacted various domains, including healthcare, where they are increasingly used for medical question answering and literature summarization. Despite their strong linguistic capabilities, LLMs suffer from a well-known limitation: hallucination, where the model generates plausible but factually incorrect or unsupported information.

This issue is particularly critical in the medical domain. Many users rely on LLMs to obtain health-related information, sometimes avoiding professional medical consultation. In such cases, inaccurate or misleading responses can lead to serious consequences, including misinterpretation of symptoms or neglect of potentially severe conditions.

To address this challenge, Retrieval-Augmented Generation (RAG) has emerged as a promising approach. By grounding LLM responses in external knowledge sources, RAG systems aim to reduce hallucinations and improve factual consistency. However, traditional RAG methods based on vector similarity often fail to capture complex relationships between medical entities.

In this work, we explore a Graph-based RAG framework that incorporates structured relationships between medical concepts. By leveraging graph representations and optimized querying strategies, the proposed approach provides richer contextual information to the LLM, enabling more accurate and evidence-based responses.

At this stage, the primary objective of this study is to evaluate the effectiveness of such an approach rather than to introduce a fully novel contribution. Future work will focus on refining the methodology and identifying unique contributions.


\section{Dataset Construction and Processing}

\subsection{Data Sources}


In this work, we focus on the mental health domain, prioritizing reliability and factual grounding. Rather than using pre-existing benchmarks, we construct a domain-specific corpus from authoritative medical sources, including:
\begin{itemize}
    \item Mental health brochures (e.g., WHO, NIMH)
    \item Clinical reference material (e.g., DSM)
\end{itemize}

\subsection{Data Extraction and Contextual Enrichment}

Raw documents (often PDF) are parsed using LLM-powered tools (e.g., LlamaIndex) and converted to Markdown. Each document is segmented into logical sections and paragraphs using Markdown headers, not arbitrary token or character limits. For each chunk, an LLM is prompted to generate a precise context sentence (e.g., ``This text belongs to the Diagnosis sub-topic of Autism Spectrum Disorder.''). This context is prepended to the text, ensuring downstream models always retain source and topical identity.

\textbf{Example data sample:}
\begin{itemize}
    \item \textbf{Context:} ``This text belongs to the Diagnosis sub-topic of Autism Spectrum Disorder.''
    \item \textbf{Text:} ``Diagnosis in young children is often a two-stage process.''
\end{itemize}

\subsection{Preprocessing}

The pipeline applies:
\begin{itemize}
    \item \textbf{Text cleaning:} Remove formatting artifacts, normalize whitespace, and correct line breaks.
    \item \textbf{Chunking:} Split by Markdown structure (headers/paragraphs), preserving semantic units (typically 1--2 paragraphs).
    \item \textbf{Contextual enrichment:} Use an LLM to generate and prepend a context sentence to each chunk.
    \item \textbf{Entity and relation extraction:} LLM-based extraction of medical entities and their relationships from each enriched chunk.
\end{itemize}

This approach preserves clinical semantics and enables accurate downstream graph construction.

\subsection{Knowledge Graph Construction}

From the enriched data, a domain-specific Knowledge Graph (KG) is built. The process includes:
\begin{itemize}
    \item \textbf{Entity canonicalization:} Normalize names (e.g., ``CBT'' $\rightarrow$ ``COGNITIVE BEHAVIORAL THERAPY''), resolve type conflicts, and filter out generic/garbage entities.
    \item \textbf{Node types:} Disorders, Symptoms, Treatments, Medications, Demographics, Safety Resources.
    \item \textbf{Relation types:} HAS\_SYMPTOM, TREATED\_BY, PRESCRIBES, SUITABLE\_FOR, CONTRAINDICATED\_FOR, URGENT\_ACTION.
    \item \textbf{LLM-based extraction:} Entities and relations are extracted using a structured LLM prompt, with pronoun resolution and strict schema validation.
    \item \textbf{Community detection:} Graph clustering (Leiden algorithm) is used to identify communities of related entities, which are then summarized by the LLM.
\end{itemize}

All steps are automated and reproducible, with checkpoints for each major artifact (enriched nodes, entities, relations, summaries).

\section{Problem Formulation}

\subsection{Task Definition}

Let $q$ denote a user query related to mental health (e.g., symptoms, diagnosis, or treatment). The objective is to generate an answer $A$ that is accurate and strictly grounded in the constructed knowledge graph.

Formally:
\begin{equation}
A = f(q, \mathcal{G})
\end{equation}
where $\mathcal{G}$ is the knowledge graph built from the enriched and extracted data.

\subsection{Retrieval Mechanism}

\textbf{Strict Graph-Only Retrieval:} Given a query $q$, the system retrieves relevant information exclusively from the knowledge graph. There is no direct text chunk retrieval or embedding-based search over raw text. Instead:
\begin{itemize}
    \item The query is mapped to relevant entities and subgraphs (paths, communities) using graph traversal and keyword/entity matching.
    \item Community detection and LLM-based summarization provide high-level context blocks for RAG.
    \item A hybrid query engine combines local (triplet-level) and global (community-level) graph retrieval, with LLM synthesis and explicit source citation.
\end{itemize}

Relevance is measured by graph-based scores (e.g., path length, node connectivity, keyword overlap in community summaries), not cosine similarity over text embeddings.

\subsection{Hallucination Reduction Objective}

Hallucination is defined as any generated information not supported by the retrieved graph facts or summaries. The system enforces grounding by:
\begin{itemize}
    \item Requiring all facts in answers to be traceable to specific graph communities or triplets.
    \item Rejecting or flagging answers when the graph does not contain sufficient information.
    \item Providing explicit source attribution (e.g., ``According to Community 4...'').
\end{itemize}

The objective is to maximize answer accuracy while minimizing hallucination:
\begin{equation}
\max \ \text{Accuracy}(A) \quad \text{and} \quad \min \ \mathcal{H}(A)
\end{equation}

This ensures all generated responses are grounded in verified, structured medical knowledge.


\end{document}